\section{Introduction}

Metaheuristic search is increasingly applied to problems whose objective
function is a numerical simulation: a finite-element sweep, a computational
fluid dynamics run, a trajectory evaluation, or a discrete-event model of a
production line. In these settings a single evaluation may take seconds to
hours, and the number of evaluations that fit into a project is a hard
constraint. The dominant cost is therefore not arithmetic but the number of
calls to the simulator.

Population-based metaheuristics were nevertheless designed under the opposite
assumption. The grey wolf optimizer (GWO) \cite{mirjalili2014gwo} maintains a
population of $N$ wolves and evaluates every one of them in every generation.
GWO's search direction, however, is produced by three positions only: the alpha,
beta and delta wolves that define the apex of the hierarchy. The remaining
wolves are attracted to the triangle these three span but do not shape it. Once
the population has clustered into a promising basin, consecutive wolves are
spatially close and, by the local smoothness of continuous objectives, their
fitness values are close as well; evaluating all of them again is largely
redundant.

The natural question is thus not how to design yet another GWO variant, but

\begin{quote}
\emph{when true evaluations are the scarce resource and candidate generation is
cheap, which candidates should receive them?}
\end{quote}

This re-frames the problem from \emph{search operator design} to
\emph{evaluation resource allocation}. The distinction matters for novelty: the
literature already contains many GWO variants that adjust the search equations,
select the leaders, or schedule operators with reinforcement learning
\cite{...}. Far fewer ask how a GWO should spend its evaluation budget.

Sparse-evaluation frameworks have recently been developed for particle swarm
optimization (PSO) and competitive swarm optimization (CSO), where a fixed
fraction of the population is evaluated per generation and the saved budget
funds additional generations \cite{ssl-pso,ss-rlpso,saq-cso}. Transferring that
idea to GWO is not immediate, for two reasons that this paper addresses.

\emph{First, GWO has no memory.} In PSO each particle carries a personal best
that was evaluated at some point in the past; in CSO the competition between
pairs creates the same persistence. Sparse evaluation in those algorithms
therefore removes evaluations while leaving the memory that guides the search
intact. GWO has no such anchor: if unevaluated wolves are allowed to act as
alpha, beta and delta, the search is steered by positions that were never
verified and may be arbitrarily bad. We solve this by storing a
\emph{verified} position separately from the virtual search position and
restricting the leader hierarchy to verified positions.

\emph{Second, in a leader-driven algorithm diversity at the crowd level is
useless.} Sparse-evaluation PSO and CSO frameworks spend part of the quota on
individuals that are far from the population centroid, arguing that they explore
untested regions. In GWO such an individual influences the search only if it
enters the elite triangle, which in a bounded high-dimensional search space
essentially never happens for an extreme outlier. Our screening rule therefore
replaces crowd-level diversity with \emph{leader-level} diversity: the
exploration quota is spent near the minority leader, i.e. on the member of the
current elite triangle that lies farthest from the leader consensus, so that a
second and third basin are kept alive rather than abandoned.

Based on these two observations we propose the Resource-Aware Sparse-Evaluation
Grey Wolf Optimizer (RASE-GWO). The virtual population is $\kappa$ times larger
than the quota, the quota is allocated by a geometry-based, surrogate-free
screening rule, and a six-state, three-action Q-learning agent adapts the split
between refinement and minority-leader probing online. The method changes none
of GWO's search equations.

The contributions are as follows.
\begin{enumerate}
  \item \textbf{A decoupled evaluation framework for leader-driven
  metaheuristics.} Virtual and verified positions are separated, and the elite
  triangle is drawn exclusively from verified positions. This makes sparse
  evaluation sound for GWO and is the mechanism that a direct transfer from
  sparse PSO/CSO frameworks lacks.
  \item \textbf{A dual-potential geometric screening rule.} Consensus
  refinement selects candidates near the leader-consensus centre; minority
  leader probing selects candidates near the leader that is farthest from that
  centre. Both scores cost $O(MD)$, no surrogate is trained, and therefore no
  approximation error is introduced.
  \item \textbf{Over-sampling as the actual resource lever.} Holding both the
  true-evaluation budget and the number of generations fixed, enlarging the
  candidate pool from $N$ to $M$ wolves while keeping the per-round quota at
  $K$ improves accuracy monotonically with the over-sampling factor. The added
  cost is arithmetic, not evaluative.
  \item \textbf{Evidence about what does not work.} We report that (i) for GWO
  the sparse-evaluation schedule of reducing the per-round quota without
  over-sampling is clearly harmful at an equal number of generations,
  (ii) crowd-level diversity scoring converts very few evaluations into
  improvements, and (iii) simply shrinking the population is clearly inferior.
  These negative results delimit the framework instead of being tuned away.
\end{enumerate}

The remainder of the paper is organised as follows.
Section~\ref{sec:related} reviews related work.
Section~\ref{sec:method} presents RASE-GWO.
Sections~\ref{sec:exp} and~\ref{sec:ablation} report the benchmark study and the
ablation.
Section~\ref{sec:conclusion} concludes.
